Figure~\ref{fig:Results_SR} shows the \met or \meff\ distributions in data
after the signal region selections (\textcolor{red}{except that on the variable shown. CHECK when I will have the plot}). 
The expected contributions from the SM backgrounds are also shown. 

\begin{figure}[p]
\centering
\begin{subfigure}[t]{0.49\textwidth}\includegraphics[width=\textwidth]{Null}\caption{}\label{fig:Results_SR3L1}\end{subfigure}
\begin{subfigure}[t]{0.49\textwidth}\includegraphics[width=\textwidth]{Null}\caption{}\label{fig:Results_SR3L2}\end{subfigure}
\begin{subfigure}[t]{0.49\textwidth}\includegraphics[width=\textwidth]{Null}\caption{}\label{fig:Results_SR0b1}\end{subfigure}
\begin{subfigure}[t]{0.49\textwidth}\includegraphics[width=\textwidth]{Null}\caption{}\label{fig:Results_SR0b2}\end{subfigure}
\begin{subfigure}[t]{0.49\textwidth}\includegraphics[width=\textwidth]{Null}\caption{}\label{fig:Results_SR1b}\end{subfigure}
\begin{subfigure}[t]{0.49\textwidth}\includegraphics[width=\textwidth]{Null}\caption{}\label{fig:Results_SR3b}\end{subfigure}
\caption{Missing transverse momentum distributions after (a)\textcolor{red}{UPDATE}, (b)\textcolor{red}{UPDATE}, 
(c) \textcolor{red}{UPDATE}, (d) \textcolor{red}{UPDATE}, (e) \textcolor{red}{UPDATE} and (f) \textcolor{red}{UPDATE} 
selection, except the \met\ requirement. The results in the signal regions are shown in the last (inclusive) bin of each plot. 
The statistical uncertainties on the background prediction are included in the uncertainty band, 
as well as the full systematic uncertainties for backgrounds with fake or non-prompt leptons, or charge-flip. 
\textcolor{red}{[UPDATE: Choose the illustrative SR plots (3 RPC, 3 RPV) and put the other in back-up]}
}
\label{fig:Results_SR} 
\end{figure} 

%For illustration, a typical SUSY signal distribution corresponding to the most relevant benchmark scenario 
%in each SR is displayed.
The detailed yields for data and the different sources of SM background in the signal regions 
are presented in Table~\ref{tab:SR_yields}. 
%None of the observed events in data contain three leptons of equal charge. 
The uncertainties amount to \textcolor{red}{[UPDATE:25--50\%]} of the total background depending on the signal region. 
The contributions listed in the ``Rare'' category in the tables is dominated by $tWZ$ and triboson production in SRs with 3 leptons,
by $\ttbar H$ and $tWZ$ production in SRs with a $b$-jet veto, 
by $\ttbar H$ and $\ttbar \ttbar$ production in SRs requiring one $b$-jet 
and by $\ttbar \ttbar$  production in SRs with three $b$-jets required.
%In all four SRs the number of data events exceeds the expectation but is consistent within the uncertainties, 
%the smallest $p$-value for the SM-only hypothesis being 0.04 for SR0b5j. 
%Out of the 14 events in the SRs, 2 of the events in SR1b and the 3 events in SR0b3j contain three leptons. 
%None of those events contain three leptons of equal charge. 

\begin{table}
\scriptsize
%\small
\begin{center}
\vspace*{-0.035\textwidth}
\begin{tabular}{|l|r|r|r|r|r|r|}
\hline
Signal Region & {\bf Rpc2L3bS } & {\bf Rpc2L3bH } & {\bf Rpc2Lsoft1b  } & {\bf Rpc2Lsoft2b } & {\bf Rpc2L0bS } & {\bf Rpc2L0bH }\\
\hline
\hline
$t\bar{t}Z\gamma^*$      &        &        &        &        &    &    \\
$t\bar{t}W$              &        &        &        &        &    &    \\
$t\bar{t}H$              &        &        &        &        &    &    \\
$t\bar{t}t\bar{t}$       &        &        &        &        &    &    \\
Diboson                  &        &        &        &        &    &    \\
Rare                     &        &        &        &        &    &    \\
Fake/non-prompt leptons  &        &        &        &        &    &    \\
Charge-flip              &        &        &        &        &    &    \\	
\hline
Total Fitted bkg         &        &        &        &        &    &    \\
\hline
Observed                 &        &        &        &        &    &    \\
\hline
$S_{\rm obs}^{95}$       &	  &	   &	    &	     &    &    \\
$S_{\rm exp}^{95}$       &	  &	   &	    &	     &    &    \\
$\sigma_{\rm{vis}}$ [fb] &        &        &        &        &    &    \\
$p_{0}$ ($\rm Z$)        &        &        &        &        &    &    \\
\hline 
\end{tabular}

\vspace*{1cm}

\begin{tabular}{|l|r|r|r|r|r|r|r|}
\hline
Signal Region & {\bf Rpc3L0bS } & {\bf Rpc3L0bH } & {\bf Rpc3L1bS } & {\bf Rpc3L1bH } & {\bf Rpc2L1bS } & {\bf Rpc2L1bH } & {\bf Rpc3LSS1b }\\
\hline
\hline
$t\bar{t}Z\gamma^*$      &        &        &        &        &    &    &     \\
$t\bar{t}W$              &        &        &        &        &    &    &     \\
$t\bar{t}H$              &        &        &        &        &    &    &     \\
$t\bar{t}t\bar{t}$       &        &        &        &        &    &    &     \\
Diboson                  &        &        &        &        &    &    &     \\
Rare                     &        &        &        &        &    &    &     \\
Fake/non-prompt leptons  &        &        &        &        &   &     &     \\
Charge-flip              &        &        &        &        &   &     &     \\	
\hline
Total Fitted bkg         &        &        &        &        &   &      &    \\
\hline
Observed                 &        &        &        &        &   &      &    \\
\hline
$S_{\rm obs}^{95}$       &	  &	   &	    &	     &   &    &      \\
$S_{\rm exp}^{95}$       &	  &	   &	    &	     &   &    &      \\
$\sigma_{\rm{vis}}$ [fb] &        &        &        &        &   &    &      \\
$p_{0}$ ($\rm Z$)        &        &        &        &        &   &    &      \\
\hline 
\end{tabular}

\vspace*{1cm}

\begin{tabular}{|l|r|r|r|r|r|r|}
\hline
Signal Region & {\bf Rpv2L1bH } & {\bf Rpv2L0b } & {\bf Rpv2L3bH } & {\bf Rpv2L3bS } & {\bf Rpv2L1bS }  & {\bf Rpv2L1bM  } \\
\hline
\hline
$t\bar{t}Z\gamma^*$      &        &        &        &        &   &        \\
$t\bar{t}W$              &        &        &        &        &   &        \\
$t\bar{t}H$              &        &        &        &        &   &        \\
$t\bar{t}t\bar{t}$       &        &        &        &        &   &        \\
Diboson                  &        &        &        &        &   &        \\
Rare                     &        &        &        &        &   &        \\
Fake/non-prompt leptons  &        &        &        &        &   &        \\
Charge-flip              &        &        &        &        &   &        \\	
\hline
Total Fitted bkg         &        &        &        &        &   &        \\
\hline
Observed                 &        &        &        &        &   &        \\
\hline
$S_{\rm obs}^{95}$       &	  &	   &	    &	     &   &	 \\
$S_{\rm exp}^{95}$       &	  &	   &	    &	     &   &	 \\
$\sigma_{\rm{vis}}$ [fb] &        &        &        &        &   &       \\
$p_{0}$ ($\rm Z$)        &        &        &        &        &   &       \\
\hline 
\end{tabular}

\vspace*{1cm}

\vspace*{-0.01\textheight}\caption[p0 and UL]{Numbers of events observed in the signal regions compared with background expectations 
obtained from the fits described in the text. The ``Rare'' category contains the contributions from associated 
production $\ttbar+WW/WZ$, as well as $t+Z/WZ/\ttbar$, $H+W/Z$, and triboson production.Empty cells (indicated by a `-') 
correspond to estimates lower than $0.01$. The p-values ($p_{0}$) give the probabilities of the observations being consistent 
with the estimated backgrounds. For an observed number of events lower than expected, the p-value is truncated at 0.5. Between 
parentheses, $p$-values are also given as the number of equivalent Gaussian standard deviations (Z).
Also shown are 95\% CL upper limits on the observed number of BSM events ($S_{\rm obs}^{95}$ ), the expected number of BSM events ($S_{\rm exp}^{95}$) and
the visible cross-section ($\sigma_{\rm{vis}}$) -- see text, given the expected number of background events (and $\pm 1\sigma$ excursions of the expectation). 
\textcolor{red}{[UPDATE:Put final numbers]}}
\label{tab:SR_yields}
\end{center}
\end{table}


In the absence of any significant deviations from the SM predictions, upper limits on possible BSM contributions to the signal regions are computed, using the SUSY benchmark scenarios described in Section~\ref{sec:intro}. 
The HistFitter framework~\cite{Baak:2014wma}, which utilises a profile-likelihood-ratio test~\cite{Cowan:2010js}, 
is used to establish 95\% confidence intervals using the CL$_\mathrm{s}$ prescription~\cite{Read_CLs}. 
The likelihood is built as the product of a Poisson probability density function describing the observed number of events in the signal region 
and Gaussian distributions constraining the nuisance parameters 
associated with the systematic uncertainties whose widths correspond to the sizes of these uncertainties; 
Poisson distributions are used instead for MC and data control region statistical uncertainties.
Correlations of a given nuisance parameter between the different sources of backgrounds and the signal are taken into account when relevant. 
The hypothesis tests are performed independently for each of the signal regions. 

Table~\ref{tab:SR_yields} presents 95\% confidence level (CL) model-independent upper limits 
on the number of observed (expected) BSM events, $S_{\rm obs}^{95}$ ($S_{\rm exp}^{95}$), that may contribute to the signal regions. 
Normalising these by the integrated luminosity $L$ of the data sample, they can be interpreted as upper limits on the visible 
BSM cross-section ($\sigma_{\rm{vis}}$), 
defined as $\sigma_{\rm{vis}}=\sigma_{\rm{prod}}\times A \times\epsilon=S_{\rm obs}^{95}/L$ where $\sigma_{\rm{prod}}$ is the production cross-section, 
$A$ the acceptance and $\epsilon$ the reconstruction efficiency. 


\begin{figure}[p]
\centering
\begin{subfigure}[t]{0.49\textwidth}\includegraphics[width=\textwidth]{Null}\caption{}\label{fig:limits_feynman_gtt}\end{subfigure}
\begin{subfigure}[t]{0.49\textwidth}\includegraphics[width=\textwidth]{Null}\caption{}\label{fig:limits_feynman_gttOffshell}\end{subfigure}
\begin{subfigure}[t]{0.49\textwidth}\includegraphics[width=\textwidth]{Null}\caption{}\label{fig:limits_feynman_gg2WZ}\end{subfigure}
\begin{subfigure}[t]{0.49\textwidth}\includegraphics[width=\textwidth]{Null}\caption{}\label{fig:limits_feynman_gg2sl}\end{subfigure}
\begin{subfigure}[t]{0.49\textwidth}\includegraphics[width=\textwidth]{Null}\caption{}\label{fig:limits_feynman_b1b1}\end{subfigure}
\begin{subfigure}[t]{0.49\textwidth}\includegraphics[width=\textwidth]{Null}\caption{}\label{fig:limits_feynman_t1t1}\end{subfigure}
\caption{
Observed and expected exclusion limits on the \gluino, \sbottomone, \stopone and \ninoone masses 
in the context of RPC SUSY scenarios with simplified mass spectra. The signal region used to obtain the limits is specified for each scenario. 
The contours of the band around the expected limit are the $\pm$1$\sigma$ results, 
including all uncertainties except the theoretical ones on the signal cross-section. The dotted lines around the observed
limit illustrate the change in the observed limit as the nominal signal cross-section is scaled up and down
by the theoretical uncertainty. All limits are computed at 95\% CL. 
The diagonal lines indicate the kinematic limit for the decays in each specified scenario.  
Results are compared with the observed limits obtained by previous ATLAS searches~\cite{paperSS3L,Aad:2016jxj,Aad:2016eki,Aad:2015iea}. 
\textcolor{red}{[UPDATE: plots and captions]}.
}
\label{fig:Results_Limits_RPC} 
\end{figure} 

\begin{figure}[p]
\centering
\begin{subfigure}[t]{0.49\textwidth}\includegraphics[width=\textwidth]{Null}\caption{}\label{fig:limits_feynm_rpv_gl313}\end{subfigure}
\begin{subfigure}[t]{0.49\textwidth}\includegraphics[width=\textwidth]{Null}\caption{}\label{fig:limits_feynm_rpv_gl321}\end{subfigure}
\begin{subfigure}[t]{0.49\textwidth}\includegraphics[width=\textwidth]{Null}\caption{}\label{fig:limits_feynm_rpv_glprime}\end{subfigure}
\begin{subfigure}[t]{0.49\textwidth}\includegraphics[width=\textwidth]{Null}\caption{}\label{fig:limits_feynm_rpv_gl112l}\end{subfigure}
\begin{subfigure}[t]{0.49\textwidth}\includegraphics[width=\textwidth]{Null}\caption{}\label{fig:limits_feynm_rpv_sd313}\end{subfigure}
\begin{subfigure}[t]{0.49\textwidth}\includegraphics[width=\textwidth]{Null}\caption{}\label{fig:limits_feynm_rpv_sd321}\end{subfigure}
\caption{
Observed and expected exclusion limits on the \gluino, \stopone and $\tilde d_R$ masses 
in the context of RPV SUSY scenarios with simplified mass spectra 
featuring $\gluino\gluino$ or $\tilde d_R\tilde d_R$ pair production with exclusive decay modes. 
The signal region used to obtain the limits is specified for each scenario. 
The contours of the band around the expected limit are the $\pm$1$\sigma$ results, 
including all uncertainties except theoretical uncertainties on the signal cross-section. The dotted lines around the observed
limit illustrate the change in the observed limit as the nominal signal cross-section is scaled up and down
by the theoretical uncertainty. All limits are computed at 95\% CL. 
The diagonal lines indicate the kinematic limit for the decays in each specified scenario, 
or in the case of $\tilde d_R\tilde d_R$ production, the region where the alternative decay channel $\tilde d_R\to \gluino d$ is open. 
Results are compared with the observed limits obtained by previous ATLAS searches~\cite{Aad:2015iea}. 
\textcolor{red}{[UPDATE: plots and captions]}.
}
\label{fig:Results_Limits_RPV} 
\end{figure}

\begin{figure}[p]
\centering
\begin{subfigure}[t]{0.49\textwidth}\includegraphics[width=\textwidth]{Null}\caption{}\label{fig:limits_feynman_gtt}\end{subfigure}
\caption{
Observed and expected exclusion limits in the $m_{1/2}$ versus $m_H^2$ plane for the NUHM2 model described in the text.
The contours of the band around the expected limit are the $\pm$1$\sigma$ results, 
including all uncertainties except theoretical uncertainties on the signal cross-section. The dotted lines around the observed
limit illustrate the change in the observed limit as the nominal signal cross-section is scaled up and down
by the theoretical uncertainty. All limits are computed at 95\% CL. 
\textcolor{red}{[UPDATE: plots and Caption]}.
}
\label{fig:Results_Limits_NUHM2} 
\end{figure} 

Exclusion limits are also set on the masses of the superpartners involved in the SUSY benchmark scenarios considered in this analysis. 
Simplified models corresponding to a single production mode and with 100\% branching ratio to a specific decay chain are used, 
with the masses of the SUSY particles not involved in the process set to very high values. 
Figures~\ref{fig:Results_Limits_RPC}, \ref{fig:Results_Limits_RPV} and \ref{fig:Results_Limits_NUHM2} show the limits on all 
the models considered in Figure~\ref{fig:feynman}. \textcolor{red}{UPDATE: comment theplots when available. Comment on the observed and 
expect limits for all models, as well as stating which SR is used for which model } 




%on the mass of the $\ninoone$ as a function of the $\gluino$ or $\sbottomone$ mass for the RPC models in Fig.~\ref{fig:feynman}. 
%For these results, asymptotic formulas~\cite{Cowan:2010js} are used to model the probability distribution of the test statistic. 
%The new limits set by this analysis can be compared 
%with the existing limits set by other ATLAS SUSY searches~\cite{paperSS3L,Aad:2016jxj,Aad:2016eki,Aad:2015iea}. 

%Signal models featuring gluino pair production with a subsequent gluino decay via $\ninotwo$ and light sleptons\\ 
%($\gluino\to q\bar q\ninotwo\to q\bar q (\ell\slepton^*/\nu\tilde{\nu}^*)\to q\bar q(\ell^+\ell^-/\nu\nu)\ninoone$) 
%are probed using SR3L1 and SR3L2 (Fig.~\ref{fig:limits_SR3L}).
%In this simplified model, the gluinos decay into $u\bar u$, $d\bar d$, $s\bar s$ or $c\bar c$ with equal probabilities, 
%and the six types of leptons are also produced in the $\tilde\chi_2^0$ decays with equal probabilities. 
%The $\ninotwo$ mass is set to $m_{\ninotwo}=(m_{\gluino} + m_{\ninoone})/2$, 
%with the $\slepton$ and $\tilde{\nu}$ masses set to $m_{\slepton,\tilde{\nu}}=(m_{\ninotwo} + m_{\ninoone})/2$. 
%For each point of the parameter space, the signal region providing the best expected limit is used. 
%Gluino masses up to $m_{\gluino}\approx\SI{1.7}{TeV}$ for a light \ninoone 
%and \ninoone masses up to $m_{\ninoone}\approx\SI{1.1}{TeV}$ for gluinos with $m_{\gluino}\approx\SI{1.4}{TeV}$ are excluded in this scenario, improving upon the limits set in Ref.~\cite{paperSS3L} by up to \SI{400}{GeV} in $\gluino$ mass and up to \SI{300}{GeV} in $\ninoone$ mass. 

%Similarly, models with gluino pair production followed by two-step gluino decays involving $\chinoonepm$ and $\ninotwo$  
%($\gluino\to q\bar q \chinoonepm \to q\bar q W\ninotwo \to  q\bar q W Z \ninoone$) 
%are probed with SR0b1 and SR0b2 (Fig.~\ref{fig:limits_SR0b}).
%In this simplified model, the gluinos decay into $u\bar u$, $d\bar d$, $s\bar s$ or $c\bar c$ with equal probabilities. 
%The $\chinoonepm$ mass is set to $m_{\chinoonepm}=(m_{\gluino} + m_{\ninoone})/2$ and
%the $\ninotwo$ mass is set to $m_{\ninotwo}=(m_{\chinoonepm} + m_{\ninoone})/2$; 
%$W$ and $Z$ bosons produced in the decay chain are not necessarily on-shell. 
%For each point of the parameter space, the signal region providing the best expected limit is used. 
%The exclusion limits in this scenario reach $m_{\gluino}\approx\SI{1.6}{TeV}$ (for light $\ninoone$) 
%and $m_{\ninoone}\approx\SI{900}{GeV}$ (for $m_{\gluino}\approx\SI{1.3}{TeV}$), 
%extending the limits set in Ref.~\cite{Aad:2016jxj,paperSS3L} by up to \SI{200}{GeV} in $\gluino$ mass and up to \SI{350}{GeV} in $\ninoone$ mass.

%Exclusion limits in a simplified model of bottom squark production with chargino-mediated $\sbottomone\to tW^-\ninoone$ decays are 
%obtained with SR1b (Fig.~\ref{fig:limits_SR1b}).
%In this model the $\chinoonepm$ mass is set to $m_{\chinoonepm}=m_{\ninoone} + \SI{100}{GeV}$.
%The limits reach mass values of $m_{\sbottomone}\approx\SI{690}{GeV}$ for a light $\ninoone$, 
%while $m_{\ninoone}\lesssim\SI{260}{GeV}$ are also excluded for $m_{\sbottomone}\approx\SI{540}{GeV}$. 
%This extends the limits set in Ref.~\cite{paperSS3L} by up to \SI{150}{GeV} in $\sbottomone$ mass and~up to \SI{120}{GeV} in $\ninoone$ mass. 

%Finally, SR3b is used to set limits on masses in a simplified model with 
%gluino pair production and $\gluino\to t\bar t\ninoone$ decays 
%via an off-shell top squark (Fig.~\ref{fig:limits_SR3b}). 
%In the compressed region where one or more of the produced top quarks is off-shell, for $m_t+m_W<m_{\gluino} - m_{\ninoone}<2m_t$, an isotropically-distributed $\gluino\to tWb\ninoone$ decay is assumed, 
%while for $2 m_W<m_{\gluino} - m_{\ninoone}<m_t+m_W$ an isotropically-distributed $\gluino\to WbWb\ninoone$ decay is used. 
%Gluino masses of $m_{\gluino}\lesssim\SI{1.45}{TeV}$ are excluded for $m_{\ninoone}\lesssim\SI{700}{GeV}$, 
%and $\ninoone$ masses up to $m_{\ninoone}\approx\SI{850}{GeV}$ are also excluded for $m_{\gluino}\approx\SI{1.25}{TeV}$. 
%Limits set in Ref.~\cite{Aad:2016eki} (dark gray line, $m_{\gluino}\approx\SI{1.75}{TeV}$ for $m_{\ninoone}\lesssim\SI{800}{GeV}$) are already beyond those reported here, but the present analysis reaches better sensitivity in the compressed region with $m_{\gluino}\lesssim 2m_t+m_{\ninoone}$. 

%Figure~\ref{fig:Results_Limits_RPV} shows the limits 
%on the masses of the $\gluino$, $\stopone$ and $\tilde d_R$ in the context of the RPV models shown in Fig.~\ref{fig:feynman}. 
%For scenarios featuring gluino pair production, prompt decays of the $\stopone$ into light SM quarks are assumed 
%via non-zero RPV couplings $\lambda''_{321}$ (for $\stopone\to \bar s\bar d$) or $\lambda''_{331}$ (for $\stopone\to \bar b\bar d$). 
%Sensitivity to these scenarios is brought by SR1b-GG, and is similar for the two \stopone\ decay modes, 
%allowing to set exclusion limits up to $m_{\gluino}\approx\SI{1.3}{TeV}$ for $m_{\stopone}\approx\SI{600}{TeV}$, 
%or probe $m_{\stopone}$ up to~$\approx\SI{1}{TeV}$. 
%These results extend the limits set in Ref.~\cite{Aad:2015iea} by up to \SI{400}{GeV} in $\gluino$ mass. 

%Pair production of right-handed like-sign down squarks is probed by SR1b-DD and SR3b-DD, 
%for squark decays $\tilde d_R\to \bar t\bar s$ or $\tilde d_R\to \bar t\bar b$ 
%through respectively non-zero RPV couplings $\lambda''_{321}$ or $\lambda''_{331}$. 
%In these scenarios, simulated events are all generated with $m_{\gluino}=\SI{2}{TeV}$ 
%to avoid contributions beyond tree level with resonant gluinos~\cite{Beenakker:1996ch,GoncalvesNetto:2012yt}; 
%the exclusion limits are however provided as function of the $\tilde d_R$ and \gluino masses, 
%by changing the NLO cross-section used for the normalization, which depends on the gluino mass. 
%For $m_{\gluino}\sim m_{\tilde d_R}$, sensitivity to $\tilde d_R\tilde d_R$ production 
%(which does not take into account contributions from $\gluino\gluino$ or $\gluino\tilde q$ production) 
%reaches $m_{\tilde d_R}\approx\SI{700}{GeV}$ for both $\tilde d_R\to \bar t\bar s$ 
%and $\tilde d_R\to \bar t\bar b$. 

